% Options for packages loaded elsewhere
\PassOptionsToPackage{unicode}{hyperref}
\PassOptionsToPackage{hyphens}{url}
%
\documentclass[
]{article}
\usepackage{amsmath,amssymb}
\usepackage{iftex}
\ifPDFTeX
  \usepackage[T1]{fontenc}
  \usepackage[utf8]{inputenc}
  \usepackage{textcomp} % provide euro and other symbols
\else % if luatex or xetex
  \usepackage{unicode-math} % this also loads fontspec
  \defaultfontfeatures{Scale=MatchLowercase}
  \defaultfontfeatures[\rmfamily]{Ligatures=TeX,Scale=1}
\fi
\usepackage{lmodern}
\ifPDFTeX\else
  % xetex/luatex font selection
\fi
% Use upquote if available, for straight quotes in verbatim environments
\IfFileExists{upquote.sty}{\usepackage{upquote}}{}
\IfFileExists{microtype.sty}{% use microtype if available
  \usepackage[]{microtype}
  \UseMicrotypeSet[protrusion]{basicmath} % disable protrusion for tt fonts
}{}
\makeatletter
\@ifundefined{KOMAClassName}{% if non-KOMA class
  \IfFileExists{parskip.sty}{%
    \usepackage{parskip}
  }{% else
    \setlength{\parindent}{0pt}
    \setlength{\parskip}{6pt plus 2pt minus 1pt}}
}{% if KOMA class
  \KOMAoptions{parskip=half}}
\makeatother
\usepackage{xcolor}
\usepackage[margin=2.54cm]{geometry}
\usepackage{color}
\usepackage{fancyvrb}
\newcommand{\VerbBar}{|}
\newcommand{\VERB}{\Verb[commandchars=\\\{\}]}
\DefineVerbatimEnvironment{Highlighting}{Verbatim}{commandchars=\\\{\}}
% Add ',fontsize=\small' for more characters per line
\usepackage{framed}
\definecolor{shadecolor}{RGB}{248,248,248}
\newenvironment{Shaded}{\begin{snugshade}}{\end{snugshade}}
\newcommand{\AlertTok}[1]{\textcolor[rgb]{0.94,0.16,0.16}{#1}}
\newcommand{\AnnotationTok}[1]{\textcolor[rgb]{0.56,0.35,0.01}{\textbf{\textit{#1}}}}
\newcommand{\AttributeTok}[1]{\textcolor[rgb]{0.13,0.29,0.53}{#1}}
\newcommand{\BaseNTok}[1]{\textcolor[rgb]{0.00,0.00,0.81}{#1}}
\newcommand{\BuiltInTok}[1]{#1}
\newcommand{\CharTok}[1]{\textcolor[rgb]{0.31,0.60,0.02}{#1}}
\newcommand{\CommentTok}[1]{\textcolor[rgb]{0.56,0.35,0.01}{\textit{#1}}}
\newcommand{\CommentVarTok}[1]{\textcolor[rgb]{0.56,0.35,0.01}{\textbf{\textit{#1}}}}
\newcommand{\ConstantTok}[1]{\textcolor[rgb]{0.56,0.35,0.01}{#1}}
\newcommand{\ControlFlowTok}[1]{\textcolor[rgb]{0.13,0.29,0.53}{\textbf{#1}}}
\newcommand{\DataTypeTok}[1]{\textcolor[rgb]{0.13,0.29,0.53}{#1}}
\newcommand{\DecValTok}[1]{\textcolor[rgb]{0.00,0.00,0.81}{#1}}
\newcommand{\DocumentationTok}[1]{\textcolor[rgb]{0.56,0.35,0.01}{\textbf{\textit{#1}}}}
\newcommand{\ErrorTok}[1]{\textcolor[rgb]{0.64,0.00,0.00}{\textbf{#1}}}
\newcommand{\ExtensionTok}[1]{#1}
\newcommand{\FloatTok}[1]{\textcolor[rgb]{0.00,0.00,0.81}{#1}}
\newcommand{\FunctionTok}[1]{\textcolor[rgb]{0.13,0.29,0.53}{\textbf{#1}}}
\newcommand{\ImportTok}[1]{#1}
\newcommand{\InformationTok}[1]{\textcolor[rgb]{0.56,0.35,0.01}{\textbf{\textit{#1}}}}
\newcommand{\KeywordTok}[1]{\textcolor[rgb]{0.13,0.29,0.53}{\textbf{#1}}}
\newcommand{\NormalTok}[1]{#1}
\newcommand{\OperatorTok}[1]{\textcolor[rgb]{0.81,0.36,0.00}{\textbf{#1}}}
\newcommand{\OtherTok}[1]{\textcolor[rgb]{0.56,0.35,0.01}{#1}}
\newcommand{\PreprocessorTok}[1]{\textcolor[rgb]{0.56,0.35,0.01}{\textit{#1}}}
\newcommand{\RegionMarkerTok}[1]{#1}
\newcommand{\SpecialCharTok}[1]{\textcolor[rgb]{0.81,0.36,0.00}{\textbf{#1}}}
\newcommand{\SpecialStringTok}[1]{\textcolor[rgb]{0.31,0.60,0.02}{#1}}
\newcommand{\StringTok}[1]{\textcolor[rgb]{0.31,0.60,0.02}{#1}}
\newcommand{\VariableTok}[1]{\textcolor[rgb]{0.00,0.00,0.00}{#1}}
\newcommand{\VerbatimStringTok}[1]{\textcolor[rgb]{0.31,0.60,0.02}{#1}}
\newcommand{\WarningTok}[1]{\textcolor[rgb]{0.56,0.35,0.01}{\textbf{\textit{#1}}}}
\usepackage{graphicx}
\makeatletter
\def\maxwidth{\ifdim\Gin@nat@width>\linewidth\linewidth\else\Gin@nat@width\fi}
\def\maxheight{\ifdim\Gin@nat@height>\textheight\textheight\else\Gin@nat@height\fi}
\makeatother
% Scale images if necessary, so that they will not overflow the page
% margins by default, and it is still possible to overwrite the defaults
% using explicit options in \includegraphics[width, height, ...]{}
\setkeys{Gin}{width=\maxwidth,height=\maxheight,keepaspectratio}
% Set default figure placement to htbp
\makeatletter
\def\fps@figure{htbp}
\makeatother
\setlength{\emergencystretch}{3em} % prevent overfull lines
\providecommand{\tightlist}{%
  \setlength{\itemsep}{0pt}\setlength{\parskip}{0pt}}
\setcounter{secnumdepth}{-\maxdimen} % remove section numbering
\ifLuaTeX
  \usepackage{selnolig}  % disable illegal ligatures
\fi
\usepackage{bookmark}
\IfFileExists{xurl.sty}{\usepackage{xurl}}{} % add URL line breaks if available
\urlstyle{same}
\hypersetup{
  pdftitle={ENV 797 - Time Series Analysis for Energy and Environment Applications \textbar{} Spring 2025},
  pdfauthor={Saemin Kim},
  hidelinks,
  pdfcreator={LaTeX via pandoc}}

\title{ENV 797 - Time Series Analysis for Energy and Environment
Applications \textbar{} Spring 2025}
\usepackage{etoolbox}
\makeatletter
\providecommand{\subtitle}[1]{% add subtitle to \maketitle
  \apptocmd{\@title}{\par {\large #1 \par}}{}{}
}
\makeatother
\subtitle{Assignment 6 - Due date 02/27/25}
\author{Saemin Kim}
\date{}

\begin{document}
\maketitle

\subsection{Directions}\label{directions}

You should open the .rmd file corresponding to this assignment on
RStudio. The file is available on our class repository on Github.

Once you have the file open on your local machine the first thing you
will do is rename the file such that it includes your first and last
name (e.g., ``LuanaLima\_TSA\_A06\_Sp25.Rmd''). Then change ``Student
Name'' on line 4 with your name.

Then you will start working through the assignment by \textbf{creating
code and output} that answer each question. Be sure to use this
assignment document. Your report should contain the answer to each
question and any plots/tables you obtained (when applicable).

When you have completed the assignment, \textbf{Knit} the text and code
into a single PDF file. Submit this pdf using Sakai.

R packages needed for this assignment: ``ggplot2'', ``forecast'',
``tseries'' and ``sarima''. Install these packages, if you haven't done
yet. Do not forget to load them before running your script, since they
are NOT default packages.

\begin{Shaded}
\begin{Highlighting}[]
\CommentTok{\#Load/install required package here}
\FunctionTok{library}\NormalTok{(ggplot2)}
\FunctionTok{library}\NormalTok{(forecast)}
\FunctionTok{library}\NormalTok{(tseries)}
\FunctionTok{library}\NormalTok{(sarima)}
\FunctionTok{library}\NormalTok{(cowplot) }
\end{Highlighting}
\end{Shaded}

This assignment has general questions about ARIMA Models.

\subsection{Q1}\label{q1}

Describe the important characteristics of the sample autocorrelation
function (ACF) plot and the partial sample autocorrelation function
(PACF) plot for the following models:

\begin{itemize}
\tightlist
\item
  AR(2)
\end{itemize}

\begin{quote}
Answer:For an AR(2) model, the ACF typically exhibits an exponentially
decaying or damped sinusoidal (oscillatory) pattern as the lags
increase. It does not cut off but instead tails off to zero. In
contrast, the PACF shows a distinct cut-off after lag 2. This means that
while the first two lags will have significant spikes, all subsequent
lags (lag 3 and beyond) should theoretically be zero or within the
statistical insignificance range.
\end{quote}

\begin{itemize}
\tightlist
\item
  MA(1)
\end{itemize}

\begin{quote}
Answer:For an MA(1) model, the ACF shows a sharp cut-off after lag 1,
meaning only the first lag is significantly different from zero. On the
other hand, the PACF does not cut off but follows an exponential decay
or a damped sinusoidal pattern, tailing off gradually as the lags
increase.
\end{quote}

\subsection{Q2}\label{q2}

Recall that the non-seasonal ARIMA is described by three parameters
ARIMA\((p,d,q)\) where \(p\) is the order of the autoregressive
component, \(d\) is the number of times the series need to be
differenced to obtain stationarity and \(q\) is the order of the moving
average component. If we don't need to difference the series, we don't
need to specify the ``I'' part and we can use the short version, i.e.,
the ARMA\((p,q)\).

\begin{enumerate}
\def\labelenumi{(\alph{enumi})}
\tightlist
\item
  Consider three models: ARMA(1,0), ARMA(0,1) and ARMA(1,1) with
  parameters \(\phi=0.6\) and \(\theta= 0.9\). The \(\phi\) refers to
  the AR coefficient and the \(\theta\) refers to the MA coefficient.
  Use the \texttt{arima.sim()} function in R to generate \(n=100\)
  observations from each of these three models. Then, using
  \texttt{autoplot()} plot the generated series in three separate
  graphs.
\end{enumerate}

\begin{Shaded}
\begin{Highlighting}[]
\FunctionTok{set.seed}\NormalTok{(}\DecValTok{123}\NormalTok{) }
\NormalTok{n }\OtherTok{\textless{}{-}} \DecValTok{100}
\NormalTok{phi }\OtherTok{\textless{}{-}} \FloatTok{0.6}
\NormalTok{theta }\OtherTok{\textless{}{-}} \FloatTok{0.9}

\CommentTok{\# model}
\NormalTok{arma10 }\OtherTok{\textless{}{-}} \FunctionTok{arima.sim}\NormalTok{(}\AttributeTok{model =} \FunctionTok{list}\NormalTok{(}\AttributeTok{ar =}\NormalTok{ phi), }\AttributeTok{n =}\NormalTok{ n)}
\NormalTok{arma01 }\OtherTok{\textless{}{-}} \FunctionTok{arima.sim}\NormalTok{(}\AttributeTok{model =} \FunctionTok{list}\NormalTok{(}\AttributeTok{ma =}\NormalTok{ theta), }\AttributeTok{n =}\NormalTok{ n)}
\NormalTok{arma11 }\OtherTok{\textless{}{-}} \FunctionTok{arima.sim}\NormalTok{(}\AttributeTok{model =} \FunctionTok{list}\NormalTok{(}\AttributeTok{ar =}\NormalTok{ phi, }\AttributeTok{ma =}\NormalTok{ theta), }\AttributeTok{n =}\NormalTok{ n)}

\CommentTok{\# plot}
\NormalTok{p1 }\OtherTok{\textless{}{-}} \FunctionTok{autoplot}\NormalTok{(arma10) }\SpecialCharTok{+} \FunctionTok{ggtitle}\NormalTok{(}\StringTok{"ARMA(1,0)"}\NormalTok{)}
\NormalTok{p2 }\OtherTok{\textless{}{-}} \FunctionTok{autoplot}\NormalTok{(arma01) }\SpecialCharTok{+} \FunctionTok{ggtitle}\NormalTok{(}\StringTok{"ARMA(0,1)"}\NormalTok{)}
\NormalTok{p3 }\OtherTok{\textless{}{-}} \FunctionTok{autoplot}\NormalTok{(arma11) }\SpecialCharTok{+} \FunctionTok{ggtitle}\NormalTok{(}\StringTok{"ARMA(1,1)"}\NormalTok{)}

\FunctionTok{plot\_grid}\NormalTok{(p1, p2, p3, }\AttributeTok{ncol =} \DecValTok{1}\NormalTok{)}
\end{Highlighting}
\end{Shaded}

\includegraphics{TSA_A06_Sp26_files/figure-latex/unnamed-chunk-2-1.pdf}

\begin{enumerate}
\def\labelenumi{(\alph{enumi})}
\setcounter{enumi}{1}
\tightlist
\item
  Plot the sample ACF for each of these models in one window to
  facilitate comparison (Hint: use \texttt{cowplot::plot\_grid()}).
\end{enumerate}

\begin{Shaded}
\begin{Highlighting}[]
\CommentTok{\# ACF plot}
\NormalTok{acf1 }\OtherTok{\textless{}{-}} \FunctionTok{ggAcf}\NormalTok{(arma10) }\SpecialCharTok{+} \FunctionTok{ggtitle}\NormalTok{(}\StringTok{"ACF: ARMA(1,0)"}\NormalTok{)}
\NormalTok{acf2 }\OtherTok{\textless{}{-}} \FunctionTok{ggAcf}\NormalTok{(arma01) }\SpecialCharTok{+} \FunctionTok{ggtitle}\NormalTok{(}\StringTok{"ACF: ARMA(0,1)"}\NormalTok{)}
\NormalTok{acf3 }\OtherTok{\textless{}{-}} \FunctionTok{ggAcf}\NormalTok{(arma11) }\SpecialCharTok{+} \FunctionTok{ggtitle}\NormalTok{(}\StringTok{"ACF: ARMA(1,1)"}\NormalTok{)}
\FunctionTok{plot\_grid}\NormalTok{(acf1, acf2, acf3, }\AttributeTok{ncol =} \DecValTok{3}\NormalTok{)}
\end{Highlighting}
\end{Shaded}

\includegraphics{TSA_A06_Sp26_files/figure-latex/unnamed-chunk-3-1.pdf}

\begin{enumerate}
\def\labelenumi{(\alph{enumi})}
\setcounter{enumi}{2}
\tightlist
\item
  Plot the sample PACF for each of these models in one window to
  facilitate comparison.
\end{enumerate}

\begin{Shaded}
\begin{Highlighting}[]
\CommentTok{\# PACF plot}
\NormalTok{pacf1 }\OtherTok{\textless{}{-}} \FunctionTok{ggPacf}\NormalTok{(arma10) }\SpecialCharTok{+} \FunctionTok{ggtitle}\NormalTok{(}\StringTok{"PACF: ARMA(1,0)"}\NormalTok{)}
\NormalTok{pacf2 }\OtherTok{\textless{}{-}} \FunctionTok{ggPacf}\NormalTok{(arma01) }\SpecialCharTok{+} \FunctionTok{ggtitle}\NormalTok{(}\StringTok{"PACF: ARMA(0,1)"}\NormalTok{)}
\NormalTok{pacf3 }\OtherTok{\textless{}{-}} \FunctionTok{ggPacf}\NormalTok{(arma11) }\SpecialCharTok{+} \FunctionTok{ggtitle}\NormalTok{(}\StringTok{"PACF: ARMA(1,1)"}\NormalTok{)}
\FunctionTok{plot\_grid}\NormalTok{(pacf1, pacf2, pacf3, }\AttributeTok{ncol =} \DecValTok{3}\NormalTok{)}
\end{Highlighting}
\end{Shaded}

\includegraphics{TSA_A06_Sp26_files/figure-latex/unnamed-chunk-4-1.pdf}

\begin{enumerate}
\def\labelenumi{(\alph{enumi})}
\setcounter{enumi}{3}
\tightlist
\item
  Look at the ACFs and PACFs. Imagine you had these plots for a data set
  and you were asked to identify the model, i.e., is it AR, MA or ARMA
  and the order of each component. Would you be able identify them
  correctly? Explain your answer.
\end{enumerate}

\begin{quote}
Answer: With a sample size of n=100, identifying the models correctly
can be challenging due to sampling variation (noise). For the ARMA(1,0)
and ARMA(0,1) models, we might observe the expected cut-off in the PACF
and ACF, respectively, but the patterns may not be perfectly clear. For
the ARMA(1,1) model, it is particularly difficult because both the ACF
and PACF are expected to tail off, making it hard to distinguish from
higher-order AR or MA models without further diagnostic tools like
AIC/BIC.
\end{quote}

\begin{enumerate}
\def\labelenumi{(\alph{enumi})}
\setcounter{enumi}{4}
\tightlist
\item
  Compare the PACF values R computed with the values you provided for
  the lag 1 correlation coefficient, i.e., does \(\phi=0.6\) match what
  you see on PACF for ARMA(1,0), and ARMA(1,1)? Should they match?
\end{enumerate}

\begin{quote}
Answer: For the ARMA(1,0) model (which is an AR(1) process), the PACF at
lag 1 should theoretically equal the ϕ coefficient (0.6). In the
simulation, the value should be close to 0.6, although not identical due
to the small sample size (n=100). However, for the ARMA(1,1) model, the
PACF at lag 1 does not directly equal ϕ because the presence of the
Moving Average (MA) component (θ) affects the initial correlations.
Therefore, they are only expected to match closely in a pure AR(1)
process.
\end{quote}

\begin{enumerate}
\def\labelenumi{(\alph{enumi})}
\setcounter{enumi}{5}
\tightlist
\item
  Increase number of observations to \(n=1000\) and repeat parts
  (b)-(e).
\end{enumerate}

\begin{Shaded}
\begin{Highlighting}[]
\CommentTok{\# 1. Set parameters for large sample size}
\FunctionTok{set.seed}\NormalTok{(}\DecValTok{123}\NormalTok{)}
\NormalTok{n\_large }\OtherTok{\textless{}{-}} \DecValTok{1000}
\NormalTok{phi }\OtherTok{\textless{}{-}} \FloatTok{0.6}
\NormalTok{theta }\OtherTok{\textless{}{-}} \FloatTok{0.9}

\CommentTok{\# 2. Simulate the models with n=1000}
\NormalTok{arma10\_large }\OtherTok{\textless{}{-}} \FunctionTok{arima.sim}\NormalTok{(}\AttributeTok{model =} \FunctionTok{list}\NormalTok{(}\AttributeTok{ar =}\NormalTok{ phi), }\AttributeTok{n =}\NormalTok{ n\_large)}
\NormalTok{arma01\_large }\OtherTok{\textless{}{-}} \FunctionTok{arima.sim}\NormalTok{(}\AttributeTok{model =} \FunctionTok{list}\NormalTok{(}\AttributeTok{ma =}\NormalTok{ theta), }\AttributeTok{n =}\NormalTok{ n\_large)}
\NormalTok{arma11\_large }\OtherTok{\textless{}{-}} \FunctionTok{arima.sim}\NormalTok{(}\AttributeTok{model =} \FunctionTok{list}\NormalTok{(}\AttributeTok{ar =}\NormalTok{ phi, }\AttributeTok{ma =}\NormalTok{ theta), }\AttributeTok{n =}\NormalTok{ n\_large)}

\CommentTok{\# 3. Plot ACF (Repeats part b)}
\NormalTok{acf1\_L }\OtherTok{\textless{}{-}} \FunctionTok{ggAcf}\NormalTok{(arma10\_large) }\SpecialCharTok{+} \FunctionTok{ggtitle}\NormalTok{(}\StringTok{"ACF: ARMA(1,0) [n=1000]"}\NormalTok{)}
\NormalTok{acf2\_L }\OtherTok{\textless{}{-}} \FunctionTok{ggAcf}\NormalTok{(arma01\_large) }\SpecialCharTok{+} \FunctionTok{ggtitle}\NormalTok{(}\StringTok{"ACF: ARMA(0,1) [n=1000]"}\NormalTok{)}
\NormalTok{acf3\_L }\OtherTok{\textless{}{-}} \FunctionTok{ggAcf}\NormalTok{(arma11\_large) }\SpecialCharTok{+} \FunctionTok{ggtitle}\NormalTok{(}\StringTok{"ACF: ARMA(1,1) [n=1000]"}\NormalTok{)}

\FunctionTok{library}\NormalTok{(cowplot)}
\FunctionTok{plot\_grid}\NormalTok{(acf1\_L, acf2\_L, acf3\_L, }\AttributeTok{ncol =} \DecValTok{3}\NormalTok{)}
\end{Highlighting}
\end{Shaded}

\includegraphics{TSA_A06_Sp26_files/figure-latex/unnamed-chunk-5-1.pdf}

\begin{Shaded}
\begin{Highlighting}[]
\CommentTok{\# 4. Plot PACF (Repeats part c)}
\NormalTok{pacf1\_L }\OtherTok{\textless{}{-}} \FunctionTok{ggPacf}\NormalTok{(arma10\_large) }\SpecialCharTok{+} \FunctionTok{ggtitle}\NormalTok{(}\StringTok{"PACF: ARMA(1,0) [n=1000]"}\NormalTok{)}
\NormalTok{pacf2\_L }\OtherTok{\textless{}{-}} \FunctionTok{ggPacf}\NormalTok{(arma01\_large) }\SpecialCharTok{+} \FunctionTok{ggtitle}\NormalTok{(}\StringTok{"PACF: ARMA(0,1) [n=1000]"}\NormalTok{)}
\NormalTok{pacf3\_L }\OtherTok{\textless{}{-}} \FunctionTok{ggPacf}\NormalTok{(arma11\_large) }\SpecialCharTok{+} \FunctionTok{ggtitle}\NormalTok{(}\StringTok{"PACF: ARMA(1,1) [n=1000]"}\NormalTok{)}

\FunctionTok{plot\_grid}\NormalTok{(pacf1\_L, pacf2\_L, pacf3\_L, }\AttributeTok{ncol =} \DecValTok{3}\NormalTok{)}
\end{Highlighting}
\end{Shaded}

\includegraphics{TSA_A06_Sp26_files/figure-latex/unnamed-chunk-5-2.pdf}

\begin{Shaded}
\begin{Highlighting}[]
\CommentTok{\# 5. Check the specific PACF value for ARMA(1,0) (Repeats part e)}
\CommentTok{\# exact value at lag 1}
\NormalTok{pacf\_result }\OtherTok{\textless{}{-}} \FunctionTok{pacf}\NormalTok{(arma10\_large, }\AttributeTok{plot =} \ConstantTok{FALSE}\NormalTok{)}
\NormalTok{val }\OtherTok{\textless{}{-}}\NormalTok{ pacf\_result}\SpecialCharTok{$}\NormalTok{acf[}\DecValTok{1}\NormalTok{]}

\FunctionTok{cat}\NormalTok{(}\StringTok{"The PACF value at lag 1 for ARMA(1,0) is:"}\NormalTok{, val, }\StringTok{"}\SpecialCharTok{\textbackslash{}n}\StringTok{"}\NormalTok{)}
\end{Highlighting}
\end{Shaded}

\begin{verbatim}
## The PACF value at lag 1 for ARMA(1,0) is: 0.5672922
\end{verbatim}

\begin{quote}
Answer: After increasing the sample size to n=1000, the ACF and PACF
plots become much cleaner and align more closely with theoretical
expectations. The ``noise'' is significantly reduced, making it much
easier to identify the cut-off points for AR(1) and MA(1). For
ARMA(1,1), the simultaneous tailing off of both plots becomes more
evident, providing a clearer diagnostic signal than the n=100 case. The
PACF at lag 1 for the ARMA(1,0) model also converges much closer to the
true ϕ=0.6.
\end{quote}

\subsection{Q3}\label{q3}

Consider the ARIMA model
\(y_t=0.7*y_{t-1}-0.25*y_{t-12}+a_t-0.1*a_{t-1}\)

\begin{enumerate}
\def\labelenumi{(\alph{enumi})}
\tightlist
\item
  Identify the model using the notation ARIMA\((p,d,q)(P,D,Q)_ s\),
  i.e., identify the integers \(p,d,q,P,D,Q,s\) (if possible) from the
  equation.
\end{enumerate}

\begin{quote}
Answer: The model is an \(ARIMA(1,0,1)(1,0,0)_{12}\).

\begin{itemize}
\tightlist
\item
  \(p=1\): The term \(0.7 y_{t-1}\) is the non-seasonal autoregressive
  (AR) component.
\item
  \(d=0\): No first-differencing is applied to the series.
\item
  \(q=1\): The term \(-0.1 a_{t-1}\) is the non-seasonal moving average
  (MA) component.
\item
  \(P=1\): The term \(-0.25 y_{t-12}\) is the seasonal autoregressive
  (SAR) component at lag 12.
\item
  \(D=0\): No seasonal differencing is applied.
\item
  \(Q=0\): There is no seasonal moving average (SMA) term.
\item
  \(s=12\): The seasonal period is 12 (monthly data).
\end{itemize}
\end{quote}

\begin{enumerate}
\def\labelenumi{(\alph{enumi})}
\setcounter{enumi}{1}
\tightlist
\item
  Also from the equation what are the values of the parameters, i.e.,
  model coefficients.
\end{enumerate}

\begin{quote}
Answer: The parameter values are as follows:

\begin{itemize}
\tightlist
\item
  \textbf{Non-seasonal AR (}\(\phi_1\)): \(0.7\)
\item
  \textbf{Non-seasonal MA (}\(\theta_1\)): \(0.1\) (Note: Using the
  notation \(a_t - \theta_1 a_{t-1}\))
\item
  \textbf{Seasonal AR (}\(\Phi_1\)): \(-0.25\)
\end{itemize}
\end{quote}

\subsection{Q4}\label{q4}

Simulate a seasonal ARIMA\((0, 1)\times(1, 0)_{12}\) model with
\(\phi =0 .8\) and \(\theta = 0.5\) using the \texttt{sim\_sarima()}
function from package \texttt{sarima}. The \(12\) after the bracket
tells you that \(s=12\), i.e., the seasonal lag is 12, suggesting
monthly data whose behavior is repeated every 12 months. You can
generate as many observations as you like. Note the Integrated part was
omitted. It means the series do not need differencing, therefore
\(d=D=0\). Plot the generated series using \texttt{autoplot()}. Does it
look seasonal?

\begin{Shaded}
\begin{Highlighting}[]
\CommentTok{\# Simulate a seasonal ARIMA(0,0,1)x(1,0,0)12 model}
\CommentTok{\# Based on instructions: d=0, D=0, non{-}seasonal MA=0.5, seasonal AR=0.8}
\FunctionTok{set.seed}\NormalTok{(}\DecValTok{456}\NormalTok{)}
\NormalTok{model\_sarima }\OtherTok{\textless{}{-}} \FunctionTok{sim\_sarima}\NormalTok{(}
  \AttributeTok{n =} \DecValTok{500}\NormalTok{, }
  \AttributeTok{model =} \FunctionTok{list}\NormalTok{(}
    \AttributeTok{ma =} \FloatTok{0.5}\NormalTok{,           }\CommentTok{\# Non{-}seasonal MA (q=1)}
    \AttributeTok{sar =} \FloatTok{0.8}\NormalTok{,          }\CommentTok{\# Seasonal AR (P=1)}
    \AttributeTok{nseasons =} \DecValTok{12}       \CommentTok{\# Seasonal period s=12}
\NormalTok{  )}
\NormalTok{)}

\CommentTok{\# Plotting}
\FunctionTok{library}\NormalTok{(ggplot2)}
\FunctionTok{autoplot}\NormalTok{(}\FunctionTok{ts}\NormalTok{(model\_sarima, }\AttributeTok{frequency =} \DecValTok{12}\NormalTok{)) }\SpecialCharTok{+} 
  \FunctionTok{ggtitle}\NormalTok{(}\StringTok{"Simulated Seasonal ARIMA(0,0,1)(1,0,0)[12]"}\NormalTok{)}
\end{Highlighting}
\end{Shaded}

\includegraphics{TSA_A06_Sp26_files/figure-latex/unnamed-chunk-6-1.pdf}

\subsection{Q5}\label{q5}

Plot ACF and PACF of the simulated series in Q4. Comment if the plots
are well representing the model you simulated, i.e., would you be able
to identify the order of both non-seasonal and seasonal components from
the plots? Explain.

\begin{Shaded}
\begin{Highlighting}[]
\NormalTok{p\_acf }\OtherTok{\textless{}{-}} \FunctionTok{ggAcf}\NormalTok{(model\_sarima, }\AttributeTok{lag.max =} \DecValTok{48}\NormalTok{) }\SpecialCharTok{+} \FunctionTok{ggtitle}\NormalTok{(}\StringTok{"ACF of Seasonal Model"}\NormalTok{)}
\NormalTok{p\_pacf }\OtherTok{\textless{}{-}} \FunctionTok{ggPacf}\NormalTok{(model\_sarima, }\AttributeTok{lag.max =} \DecValTok{48}\NormalTok{) }\SpecialCharTok{+} \FunctionTok{ggtitle}\NormalTok{(}\StringTok{"PACF of Seasonal Model"}\NormalTok{)}
\FunctionTok{plot\_grid}\NormalTok{(p\_acf, p\_pacf, }\AttributeTok{nrow =} \DecValTok{1}\NormalTok{)}
\end{Highlighting}
\end{Shaded}

\includegraphics{TSA_A06_Sp26_files/figure-latex/unnamed-chunk-7-1.pdf}

\begin{quote}
Answer: Yes, the plots represent the simulated model well, and the key
components can be identified as follows:

\begin{itemize}
\tightlist
\item
  \textbf{Non-seasonal components:} In the \textbf{ACF}, there is a
  significant spike at lag 1 followed by a sharp drop, which correctly
  indicates an \(MA(1)\) process (\(q=1\)). In the \textbf{PACF}, the
  values tail off for the first few lags, which is consistent with an MA
  model.
\item
  \textbf{Seasonal components (}\(s=12\)): In the \textbf{ACF}, we
  observe significant spikes at seasonal lags (12, 24, 36, etc.) that
  gradually decay (tail off), suggesting a seasonal autoregressive
  process. In the \textbf{PACF}, there is a prominent spike at lag 12
  followed by a cut-off at subsequent seasonal lags (24, 36, etc.). This
  clearly identifies a \(SAR(1)\) process (\(P=1\)).
\end{itemize}

Overall, the combination of a cut-off in the ACF for the non-seasonal
part and a cut-off in the PACF for the seasonal part allows us to
correctly identify the simulated \(ARIMA(0,0,1) \times (1,0,0)_{12}\)
model.
\end{quote}

\end{document}
